\subsubsection*{\texttt{34-retrieval-topk5}}
\label{app:samples:34_retrieval_topk5}

\paragraph{Q3: How long does the Haste spell last?}
\noindent \textbf{Corr.:}~2.00/5\quad \textbf{Faith.:}~0.00\quad \textbf{Relev.:}~0.00\quad \textbf{Ctx.Rel.:}~1.00\\

\noindent\textbf{Antwort:} The Haste spell lasts for the duration of concentration up to 1 minute. At higher levels, the duration can be increased to 1 year with a 7th or 8th level spell slot, or indefinitely with a 9th level spell slot.

\noindent\textbf{Abgerufene Chunks:}
\begin{itemize}[noitemsep, topsep=2pt]
  \item[\textbf{\#1}] \texttt{[0.0167]} \textit{Spells/Haste} -- \#\#\# Haste  *3rd-level transmutation*  **Casting Time:** 1 action  **Range:** 30 \ldots
  \item[\textbf{\#2}] \texttt{[0.0164]} \textit{Spells/# Spellcasting} -- \#\# Duration  A spell's duration is the length of time the spell persists. A dura\ldots
  \item[\textbf{\#3}] \texttt{[0.0161]} \textit{Spells/Modify Memory} -- \#\#\# Modify Memory  *5th-level enchantment*  **Casting Time:** 1 action  **Range:\ldots
  \item[\textbf{\#4}] \texttt{[0.0159]} \textit{Spells/Geas} -- \#\#\# Geas  *5th-level enchantment*  **Casting Time:** 1 minute  **Range:** 60 fee\ldots
  \item[\textbf{\#5}] \texttt{[0.0156]} \textit{Treasure/Tome of Understanding} -- \#\#\# Tome of Understanding  *Wondrous item, very rare*  This book contains intuit\ldots
\end{itemize}
\vspace{4pt}
\hrule
\vspace{6pt}

\paragraph{Q5: Can a grappled creature cast spells?}
\noindent \textbf{Corr.:}~2.00/5\quad \textbf{Faith.:}~1.00\quad \textbf{Relev.:}~0.00\quad \textbf{Ctx.Rel.:}~0.25\\

\noindent\textbf{Antwort:} The provided rules do not specify whether a grappled creature can cast spells or not. Therefore, I cannot determine the answer from the context alone.

\noindent\textbf{Abgerufene Chunks:}
\begin{itemize}[noitemsep, topsep=2pt]
  \item[\textbf{\#1}] \texttt{[0.0167]} \textit{Gameplay/Combat} -- \#\# Cast a Spell  Spellcasters such as wizards and clerics, as well as many monst\ldots
  \item[\textbf{\#2}] \texttt{[0.0164]} \textit{Monsters/# Monster Statistics} -- \#\#\# Innate Spellcasting  A monster with the innate ability to cast spells has th\ldots
  \item[\textbf{\#3}] \texttt{[0.0161]} \textit{Classes/Sorcerer} -- \#\#\#\# Careful Spell  When you cast a spell that forces other creatures to make a \ldots
  \item[\textbf{\#4}] \texttt{[0.0159]} \textit{Gameplay/Combat} -- \#\#\# Grappling  When you want to grab a creature or wrestle with it, you can use \ldots
  \item[\textbf{\#5}] \texttt{[0.0156]} \textit{Gamemastering/Conditions} -- \#\# Grappled  - A grappled creature's speed becomes 0, and it can't benefit from \ldots
\end{itemize}
\vspace{4pt}
\hrule
\vspace{6pt}

\paragraph{Q12: If a creature has exhaustion level 3 and is also restrained, what is its movement speed?}
\noindent \textbf{Corr.:}~5.00/5\quad \textbf{Faith.:}~1.00\quad \textbf{Relev.:}~1.00\quad \textbf{Ctx.Rel.:}~0.50\\

\noindent\textbf{Antwort:} The creature's movement speed is disabled due to the restrained condition. Therefore, its speed is 0.

\noindent\textbf{Abgerufene Chunks:}
\begin{itemize}[noitemsep, topsep=2pt]
  \item[\textbf{\#1}] \texttt{[0.0167]} \textit{Gamemastering/Conditions} -- \#\# Restrained  - A restrained creature's speed becomes 0, and it can't benefit f\ldots
  \item[\textbf{\#2}] \texttt{[0.0164]} \textit{Gameplay/Adventuring} -- \#\# Suffocating  A creature can hold its breath for a number of minutes equal to \ldots
  \item[\textbf{\#3}] \texttt{[0.0161]} \textit{Gamemastering/Conditions} -- \#\# Exhaustion  Some special abilities and environmental hazards, such as starvat\ldots
  \item[\textbf{\#4}] \texttt{[0.0159]} \textit{Spells/Freedom of Movement} -- \#\#\# Freedom of Movement  *4th-level abjuration*  **Casting Time:** 1 action  **R\ldots
  \item[\textbf{\#5}] \texttt{[0.0156]} \textit{Gamemastering/Conditions} -- \#\# Unconscious  - An unconscious creature is incapacitated (see the condition), \ldots
\end{itemize}
\vspace{4pt}
\hrule
\vspace{6pt}

\paragraph{Q16: Can you use Action Surge to cast two leveled spells in one turn?}
\noindent \textbf{Corr.:}~2.00/5\quad \textbf{Faith.:}~1.00\quad \textbf{Relev.:}~1.00\quad \textbf{Ctx.Rel.:}~0.25\\

\noindent\textbf{Antwort:} Based on the provided context, no, you cannot use Action Surge to cast two leveled spells in one turn. Action Surge allows you to take an additional action on top of your regular action, but casting a spell still requires a specific type of action (bonus action or regular action) as indicated in the\ldots

\noindent\textbf{Abgerufene Chunks:}
\begin{itemize}[noitemsep, topsep=2pt]
  \item[\textbf{\#1}] \texttt{[0.0167]} \textit{Classes/Fighter} -- \#\#\# Action Surge  Starting at 2nd level, you can push yourself beyond your norma\ldots
  \item[\textbf{\#2}] \texttt{[0.0164]} \textit{Spells/# Spellcasting} -- \#\#\# Bonus Action  A spell cast with a bonus action is especially swift. You must\ldots
  \item[\textbf{\#3}] \texttt{[0.0161]} \textit{Classes/Sorcerer} -- \#\#\#\# Quickened Spell  When you cast a spell that has a casting time of 1 action,\ldots
  \item[\textbf{\#4}] \texttt{[0.0159]} \textit{Characterizations/Multiclassing} -- If you have more than one spellcasting class, this table might give you spell sl\ldots
  \item[\textbf{\#5}] \texttt{[0.0156]} \textit{Classes/Sorcerer} -- \#\#\#\# Flexible Casting  You can use your sorcery points to gain additional spell \ldots
\end{itemize}
\vspace{4pt}
\hrule
\vspace{6pt}
